\documentclass[10pt,a4paper]{amsart}
\usepackage[margin=19mm]{geometry}
\usepackage{amsmath,amssymb,mathtools}
\usepackage[scr=rsfs,cal=euler]{mathalfa}
\usepackage{xfrac}
\usepackage[unicode,hidelinks,bookmarks=false]{hyperref}
\newcommand{\ie}{i.e.\ }
\newcommand{\cf}{cf.\ }
\newcommand{\resp}{respectively\ }
\newcommand{\Subsection}{\subsection}
\providecommand{\nobreakdash}{\nobreak\hbox{-}}
\setlength{\emergencystretch}{2em}
\multlinegap0pt
\usepackage{amssymb}
\newtheorem{lemma}{Lemma}
 \newcommand{\Mod}[1]{\ (\mathrm{mod}\ #1)}
\newcommand{\fr}{\frac}
\title{Congruences for two-color partitions with odd smallest part}
\author{George E. Andrews, Mohamed El Bachraoui}
\date{Source: arXiv:2410.14190v2 (CC BY 4.0)}
\begin{document}
\maketitle

\noindent\emph{Source and licence.} Congruences for two-color partitions with odd smallest part. Version arXiv:2410.14190v2 (CC BY 4.0), distributed by arXiv under the Creative Commons Attribution 4.0 International licence, http://creativecommons.org/licenses/by/4.0/, as recorded in the arXiv licence metadata for this exact version and shown on the abstract page https://arxiv.org/abs/2410.14190v2. Full licence notice: \texttt{licenses/arxiv-2410.14190-CC-BY-4.0.txt}.\par\smallskip

\noindent\textit{Editorial scope.} Counted unit: the article's lemma cong3 (source0.tex lines 626--636). The article's complete original proof of it is delivered with it (source0.tex lines 638--679, one \texttt{proof} environment of 42 lines). No mathematics is written, repaired or re-derived: the statement and the proof are the article's own lines, in the article's own notation, and no retained line is substituted or re-flowed.  Retained uncounted as the source-local context the proof needs: lines 621--624.  Nothing was omitted from any retained span, so the package carries no omission record.  No retained line contains a citation, so the wrapper carries no bibliography and no bibliography entry.  No retained line contains a figure environment, an image-inclusion command or drawing code, so no image is delivered and the package declares no asset dependency.  Editorial formatting only, in addition: an AMS \texttt{amsart} preamble that restates the article's own declarations and macros used by the retained lines, namely the environments \texttt{lemma} on separate counters, together with the standard packages those lines need.  The retained environments print in this document as Lemma 1 in order of appearance, whereas the article prints the same items with the numbering of its own full document.  The complete native source of the article as distributed in the arXiv e-print for this exact version is bundled unchanged as \texttt{sources/167037/source0.tex}.
\begin{equation}\label{basic-cong mod2}
1+q \equiv 1-q \Mod{2}\ \text{and\ }
(1+q)^2 \equiv 1+q^2\Mod{2},
\end{equation}

\begin{lemma}\label{lem cong3}
Let
\[
f(q)=\frac{q^5(1-q^2)^2(1-q^4)}{(1+q^2)(1+q^4)(1-q)^2(1-q^3)^2}\ \text{and\ }
g(q)=\frac{q^5(1-q^4)}{(1+q^2)(1-q)^2(1-q^3)^2}.
\]
Then for all integers $n\ge0$ we have
\[
[q^{4n}]f(q)\equiv [q^{4n}]g(q)\Mod{4}.
\]
\end{lemma}

\begin{proof}
We compute
\[
\begin{aligned}
f(q)-g(q)
&=\frac{q^5(1-q^4)\bigl((1-q^2)^2-(1+q^4)\bigr)}
        {(1+q^2)(1+q^4)(1-q)^2(1-q^3)^2} \\
&=\frac{-2q^7(1-q^4)}{(1+q^2)(1+q^4)(1-q)^2(1-q^3)^2} \\
&=-2q^7 h(q),
\end{aligned}
\]
where
\[
h(q)=\fr{(1-q^4)}{(1+q^2)(1+q^4)(1-q)^2(1-q^3)^2} := \sum_{n\ge0} c_n q^n.
\]
Then
\[
[q^{4n}]\big( f(q)-g(q) \big)=-2\,c_{4n-7}.
\]
Thus to prove $[q^{4n}]f(q)\equiv [q^{4n}]g(q)\Mod{4}$, it suffices to show that
$c_{4n-7}$ is even for all $n$.
With the help of~\eqref{basic-cong mod2}, we get
\begin{align*}
h(q) &\equiv \frac{1+q^4}{(1+q^2)(1+q^4)(1+q)^2(1+q^3)^2} \Mod{2} \\
&\equiv \frac{1}{(1+q^2)(1+q)^2(1+q^3)^2}\Mod{2} \\
&\equiv \frac{1}{(1+q^2)(1+q^2)(1+q^6)} \Mod{2}\\
&\equiv \frac{1}{(1+q^2)^2(1+q^6)}\Mod{2} \\
&\equiv \Big(\sum_{k\ge0} q^{2k}\Big)^2 \Big(\sum_{l\ge0} q^{6l}\Big) \Mod{2} \\
&\equiv \Big(\sum_{k\ge0} q^{4k}\Big)\Big(\sum_{l\ge0} q^{6l}\Big) \\
&=\sum_{k,l\geq 0} q^{4k+6l}.
\end{align*}
%
Consequently, for any integer $n\geq 0$,
\[
[q^{2n+1}]h(q)\equiv 0 \Mod{2}
\]
and therefore
\[
c_{4n-7} \equiv [q^{4n-7}]h(q)\equiv 0\Mod{2}.
\]
This completes the proof.
\end{proof}

\end{document}
